% SGA 3, Expose III, Loop 1 English reconstruction.
% Scope: Reminders 0.1.7 through Corollary 0.1.10, beginning immediately
% after Remark 0.1.6 and stopping immediately before Lemma 0.1.11.
% Authority: Polo--Gille Exp3-14oct24.pdf, component-local pages 5--12
% (combined-reader pages 115--122). Force-OCR is a drafting locator; English
% wording was compared with Jacob Reinhold's CC BY 4.0 Markdown lineage.
% The PDF controls text, formulas, numbering, notes, and all diagram images.

% Source page: Exp3 component-local p. 5; combined-reader p. 115.
\medskip
\noindent\SGAInlineRefTarget{sga3:iii:paragraph:0-1-7}\textbf{回顾 0.1.7.}
\footnote{编者注: 本段为新增内容 (另见 [DG70], \S I.4, 第~1--2~条).}
本段对给定的 \(S\)-概形 \(X\), ``回顾''微分模
\(\Omega^1_{X/S}\) 和对角形的第一无穷小邻域
\(\Delta^{(1)}_{X/S}\) 的某些函子性质. 这些性质容易由
EGA~IV\textsubscript{4}, \S\S16.1--16.4 得到, 但在那里没有明确表述.

\smallskip
\noindent\textit{(a)} 先回顾以下事实 (参见 EGA~II, \S\S1.2--1.5).
设 \(g:Y\to X\) 是概形态射, \(\pi:X'\to X\) 是仿射
\(X\)-概形, \(\mathscr B=\pi_*(\mathscr O_{X'})\) 是相应的
拟凝聚 \(\mathscr O_X\)-代数. 那么 \(Y\)-概形 \(Y\times_X X'\)
仿射, 并对应于拟凝聚 \(\mathscr O_Y\)-代数 \(g^*(\mathscr B)\),
且有由双射构成的交换图:

\SGAThreeNativeDiagram[0.82\textwidth]
  {figures/exp3/diagram_003_exp3_p005_reminders017_affine_bijections.tex}
  {沿 \(g\) 作基变换后的仿射概形与代数之间的双射.}
  {fig:exp3-reminders-017-affine-bijections}
  {Polo--Gille Exp3-14oct24.pdf, 组件第 5 页; 合卷第 115 页}

此外, 这些双射关于二元组 \((X,X')\) 具有函子性. 更确切地说,
设 \(W'\) 是对应于拟凝聚 \(\mathscr O_W\)-代数 \(\mathscr A\)
的仿射 \(W\)-概形, 并假设有以下概形态射构成的交换图.

\SGAThreeNativeDiagram[0.38\textwidth]
  {figures/exp3/diagram_004_exp3_p005_reminders017_scheme_functoriality.tex}
  {用以表达关于 \((X,X')\) 的函子性的概形态射.}
  {fig:exp3-reminders-017-scheme-functoriality}
  {Polo--Gille Exp3-14oct24.pdf, 组件第 5 页; 合卷第 115 页}

将与 \(f'\) 关联的代数态射记为
\[
 \varphi:\mathscr A\longrightarrow f_*(\mathscr B),
 \qquad
 \varphi^\sharp:f^*(\mathscr A)\longrightarrow\mathscr B
\]
并将与变动的 \(X\)-态射 \(g':Y\to X'\) 关联的代数态射记为
\[
 \theta:\mathscr B\longrightarrow g_*(\mathscr O_Y),
 \qquad
 \theta^\sharp:g^*(\mathscr B)\longrightarrow\mathscr O_Y
\]
则下图交换, 其中 \(Y\) 经由 \(f\circ g\) 视为 \(W\)-概形.

% Source page: Exp3 component-local p. 6; combined-reader p. 116.
\SGAThreeNativeDiagram[0.96\textwidth]
  {figures/exp3/diagram_005_exp3_p006_reminders017_hom_functoriality.tex}
  {仿射概形与代数之间的对应的函子性.}
  {fig:exp3-reminders-017-hom-functoriality}
  {Polo--Gille Exp3-14oct24.pdf, 组件第 6 页; 合卷第 116 页}

\smallskip
\noindent\SGAInlineRefTarget{sga3:iii:paragraph:0-1-7:item:b}\textit{(b)}
现设 \(X\) 是 \(S\)-概形. 令 \(\Omega^1_{X/S}\) 为 \(X\) 在
\(S\) 上的微分模, \(\Delta^{(1)}_{X/S}\) 为对角浸入
\(\delta_X:X\to X\times_S X\) 的第一无穷小邻域.
它是 \(X\times_S X\) 的一个子概形, 对角形 \(\Delta_{X/S}\)
是其中的闭子概形. 用 \(\operatorname{pr}^i_X\), \(i=1,2\),
表示两个投影 \(X\times_S X\to X\), 并用 \(\pi_X\) 表示
\(\operatorname{pr}^1_X\) 在 \(\Delta^{(1)}_{X/S}\) 上的限制.

一方面, 每个 \(S\)-概形态射 \(f:X\to W\) 都诱导一个
\(S\)-态射
\[
 \Delta^{(1)}f:\Delta^{(1)}_{X/S}\longrightarrow
 \Delta^{(1)}_{W/S}
\]
使下图交换.

\SGAThreeNativeDiagram[0.86\textwidth]
  {figures/exp3/diagram_006_exp3_p006_reminders017_delta_functoriality.tex}
  {对角形的第一无穷小邻域的函子性.}
  {fig:exp3-reminders-017-delta-functoriality}
  {Polo--Gille Exp3-14oct24.pdf, 组件第 6 页; 合卷第 116 页}

另一方面, 经由投影 \(\pi_X\), 概形 \(\Delta^{(1)}_{X/S}\)
是一个仿射 \(X\)-概形: 它是增广拟凝聚 \(\mathscr O_X\)-代数
\[
 \mathscr P^1_{X/S}=\mathscr O_X\oplus\Omega^1_{X/S},
\]
的谱, 其中 \(\Omega^1_{X/S}\) 是一个平方为零的理想.
增广是指在 \(\Omega^1_{X/S}\) 上为零的 \(\mathscr O_X\)-代数态射
\[
 \varepsilon_X:\mathscr P^1_{X/S}\longrightarrow\mathscr O_X
\]
它对应于闭浸入 \(\delta_X:X\hookrightarrow\Delta^{(1)}_{X/S}\).
因此, 每个 \(S\)-概形态射 \(f:X\to W\) 都诱导一个增广
\(\mathscr O_X\)-代数态射
\[
 f^*(\mathscr P^1_{W/S})
 =\mathscr O_X\oplus f^*(\Omega^1_{W/S})
 \longrightarrow
 \mathscr P^1_{X/S}=\mathscr O_X\oplus\Omega^1_{X/S},
\]
或者等价地, 诱导一个 \(\mathscr O_X\)-模态射
\[
 f_{X/W/S}:f^*(\Omega^1_{W/S})\longrightarrow\Omega^1_{X/S};
\]
关于记号 \(f_{X/W/S}\), 参见 EGA~IV\textsubscript{4},
(16.4.3.6) 和 (16.4.18.2).

% Source page: Exp3 component-local p. 7; combined-reader p. 117.
由于 \(\pi_X:\Delta^{(1)}_{X/S}\to X\) 仿射, 由 (a) 可知,
\(\Delta^{(1)}f\) 完全由 \(f_{X/W/S}\) 决定. 对每个
\(X\)-概形 \(g:Y\to X\), 集合
\[
 \operatorname{Hom}_X(Y,\Delta^{(1)}_{X/S})
 \simeq
 \operatorname{Hom}_{\mathscr O_Y\text{-}\mathrm{alg}}
 \bigl(\mathscr O_Y\oplus g^*(\Omega^1_{X/S}),\mathscr O_Y\bigr)
\]
可等同于 \(\operatorname{Hom}_{\mathscr O_Y}
 (g^*(\Omega^1_{X/S}),\mathscr O_Y)\) 的一个子集, 即
\[
 \widetilde{\operatorname{Hom}}_{\mathscr O_Y}
 \bigl(g^*(\Omega^1_{X/S}),\mathscr O_Y\bigr),
\]
它由那些满足 \(\operatorname{Im}(\psi)\) 是 \(\mathscr O_Y\)
中平方为零的理想的 \(\mathscr O_Y\)-态射
\(\psi:g^*(\Omega^1_{X/S})\to\mathscr O_Y\) 组成.
\footnote{编者注: 在 SGA~1, III~5.1 的证明中,
``但代数同态 \(\mathscr A\to\mathscr O_Y\) 对应于 \ldots''
这一句必须据此更正; 所引段落的其余部分不变.}

因此, 将 (a) 应用于下图.

\SGAThreeNativeDiagram[0.48\textwidth]
  {figures/exp3/diagram_007_exp3_p007_reminders017_delta_application.tex}
  {应用仿射函子性陈述的交换图.}
  {fig:exp3-reminders-017-delta-application}
  {Polo--Gille Exp3-14oct24.pdf, 组件第 7 页; 合卷第 117 页}

考虑到 \(\Delta^{(1)}f\) 是 \(f\times f\) 在
\(\Delta^{(1)}_{X/S}\) 上的限制, 就得到下面关于
\(X\)-概形 \(g:Y\to X\) 具有函子性的交换图.

\SGAThreeNativeDiagram[0.97\textwidth]
  {figures/exp3/diagram_008_exp3_p007_reminders017_star_diagram.tex}
  {第一邻域的态射与微分映射之间的函子性比较, 式 \(0.1.7\,(*)\).}
  {fig:exp3-reminders-017-star}
  {Polo--Gille Exp3-14oct24.pdf, 组件第 7 页; 合卷第 117 页}

\medskip
\noindent\SGAInlineRefTarget{sga3:iii:remark:0-1-7-1}\textbf{注 0.1.7.1.}
以如下后文将用到的注结束本段 (参见
\SGARefLink{sga3:iii:corollary:0-1-10}{0.1.10}). 用 \(\widetilde L_X\)
表示对每个 \(X\)-概形 \(g:Y\to X\) 赋予集合
\[
 \operatorname{Hom}_{\mathscr O_Y}
 \bigl(g^*(\Omega^1_{X/S}),\mathscr O_Y\bigr),
\]
的 \(X\)-函子, 并用 \(\widetilde L_f:\widetilde L_X\to\widetilde L_W\)
表示上面定义的函子态射: 它将 \(\psi\in\widetilde L_X(Y)\) 映到
\(\psi\circ g^*(f_{X/W/S})\). 上述讨论给出一个函子交换图.

\SGAThreeNativeDiagram[0.74\textwidth]
  {figures/exp3/diagram_009_exp3_p007_remark0171_functor_diagram.tex}
  {\(\widetilde L_X\) 与第一无穷小邻域的函子性.}
  {fig:exp3-remark-0171-functor-diagram}
  {Polo--Gille Exp3-14oct24.pdf, 组件第 7 页; 合卷第 117 页}

% Source page: Exp3 component-local p. 8; combined-reader p. 118.
\medskip
\noindent\SGAInlineRefTarget{sga3:iii:theorem:0-1-8}\textbf{定理 0.1.8.}
(SGA~1, III~5.1)\footnote{编者注: 这里重述 SGA~1, III~5.1
的陈述与证明, 因为后文证明 \SGARefLink{sga3:iii:proposition:0-2}{命题~0.2}
时将会用到它们; 也可先看下面的
\SGARefLink{sga3:iii:corollary:0-1-9}{推论~0.1.9}. 此外,
\(\mathscr J\) (分别地, \(\mathscr I\)) 是字母 \(J\)
(分别地, \(I\)) 的另一种书法体形式. 回到此前的记号, 即
\(\mathscr J\subset\mathscr I\) 是 \(\mathscr O_S\) 的理想,
满足 \(\mathscr I\mathscr J=0\), 应用这一结果时将取
\(\mathscr J=\mathscr J\mathscr O_Y\) (分别地,
\(\mathscr I=\mathscr I\mathscr O_Y\)).}
\textit{设 \(Y\) 与 \(X\) 是两个 \(S\)-概形,
\(\mathscr J\) 是 \(\mathscr O_Y\) 的一个平方为零的拟凝聚理想,
\(Y_{\mathscr J}\) 是 \(Y\) 中由 \(\mathscr J\) 定义的闭子概形,
\(g_{\mathscr J}:Y_{\mathscr J}\to X\) 是一个 \(S\)-态射.}
\begin{enumerate}
\item[\SGAInlineRefTarget{sga3:iii:theorem:0-1-8:item:a}(a)]
  \textit{将 \(g_{\mathscr J}\) 扩张的 \(S\)-态射 \(g:Y\to X\)
  组成的集合 \(P(g_{\mathscr J})\), 或者为空, 或者在下列交换群下为主齐性集:}
  \[
   \operatorname{Hom}_{\mathscr O_{Y_{\mathscr J}}}
   \bigl(g_{\mathscr J}^*(\Omega^1_{X/S}),\mathscr J\bigr).
  \]
\item[\SGAInlineRefTarget{sga3:iii:theorem:0-1-8:item:b}(b)]
  \textit{假设 \(i:Y_0\hookrightarrow Y_{\mathscr J}\) 是由拟凝聚理想
  \(\mathscr I\supset\mathscr J\) 定义的闭浸入, 其中
  \(\mathscr I\mathscr J=0\), 并令 \(g_0=g_{\mathscr J}\circ i\).
  则上述交换群同构于}
  \[
   \operatorname{Hom}_{\mathscr O_{Y_0}}
   \bigl(g_0^*(\Omega^1_{X/S}),\mathscr J\bigr).
  \]
\end{enumerate}

\noindent\textit{证明.}
\SGARefLink{sga3:iii:theorem:0-1-8:item:b}{(b)} 立即由
\SGARefLink{sga3:iii:theorem:0-1-8:item:a}{(a)} 得到. 事实上,
由于 \(\mathscr J\) 被 \(\mathscr I\) 零化, 可将其视为
\(\mathscr O_{Y_0}\)-模, 而由伴随性得
\[
 \operatorname{Hom}_{\mathscr O_{Y_{\mathscr J}}}
  \bigl(g_{\mathscr J}^*(\Omega^1_{X/S}),\mathscr J\bigr)
 =
 \operatorname{Hom}_{\mathscr O_{Y_0}}
  \bigl(i^*g_{\mathscr J}^*(\Omega^1_{X/S}),\mathscr J\bigr).
\]

为证明 \SGARefLink{sga3:iii:theorem:0-1-8:item:a}{(a)}, 可假设
\(P(g_{\mathscr J})\ne\emptyset\), 即存在扩张 \(g_{\mathscr J}\)
的 \(S\)-态射 \(g:Y\to X\). 用 \(j:Y_{\mathscr J}\hookrightarrow Y\)
表示该浸入. 则 \(P(g_{\mathscr J})\) 是满足
\(g'\circ j=g_{\mathscr J}\) 的 \(S\)-态射 \(g':Y\to X\)
组成的集合. 给定这样的 \(g'\), 等价于给定 \(S\)-态射
\[
 h:Y\longrightarrow X\times_S X
\]
使得 \(\operatorname{pr}_1\circ h=g\) 及
\(h_{\mathscr J}=\delta\circ g_{\mathscr J}\), 其中
\(h_{\mathscr J}=h\circ j\), \(\delta:X\hookrightarrow X\times_S X\)
是对角浸入.

\SGAThreeNativeDiagram[0.56\textwidth]
  {figures/exp3/diagram_010_exp3_p008_theorem018_factorization.tex}
  {由 \(g_{\mathscr J}\) 的两个扩张确定的态射 \(h\).}
  {fig:exp3-theorem-018-factorization}
  {Polo--Gille Exp3-14oct24.pdf, 组件第 8 页; 合卷第 118 页}

由于 \(h_{\mathscr J}\) 经过 \(\delta\) 分解, 又由于
\(\mathscr J^2=0\), \(Y\) 位于浸入 \(j:Y_{\mathscr J}\to Y\)
的第一无穷小邻域内, 由函子性 (EGA~IV\textsubscript{4}, 16.2.2(i))
可知, 所求态射 \(h\) 唯一地经过 \(\Delta^{(1)}_{X/S}\) 分解;
参见 0.1.7. 令
\[
 Y'=\Delta^{(1)}_{X/S}\times_XY,
 \qquad
 Y'_{\mathscr J}=Y'\times_YY_{\mathscr J}
 =\Delta^{(1)}_{X/S}\times_XY_{\mathscr J}.
\]
于是, 所求态射 \(h\) 与 \(Y'\to Y\) 的那些扩张
\(Y'_{\mathscr J}\to Y_{\mathscr J}\) 的截面
\(u_{\mathscr J}=(\delta\circ g_{\mathscr J},\operatorname{id})\)
的截面 \(u\) 一一对应. 另一方面, \(Y'\) (分别地,
\(Y'_{\mathscr J}\)) 在 \(Y\) (分别地, \(Y_{\mathscr J}\)) 上仿射,
并对应于拟凝聚代数
\[
 \mathscr A=\mathscr O_Y\oplus g^*(\Omega^1_{X/S}),
 \qquad\text{分别地}\qquad
 \mathscr A_{\mathscr J}
 =\mathscr A\otimes_{\mathscr O_Y}\mathscr O_{Y_{\mathscr J}}
 =\mathscr O_{Y_{\mathscr J}}
   \oplus g_{\mathscr J}^*(\Omega^1_{X/S}).
\]
令 \(\varepsilon:\mathscr A\to\mathscr O_Y\) 为典范增广,
即在 \(g^*(\Omega^1_{X/S})\) 上为零的 \(\mathscr O_Y\)-代数态射,
并类似地定义 \(\varepsilon_{\mathscr J}:\mathscr A_{\mathscr J}
 \to\mathscr O_{Y_{\mathscr J}}\). 则
\[
 \Gamma(Y'/Y)
 \simeq\operatorname{Hom}_{\mathscr O_Y\text{-}\mathrm{alg}}
       (\mathscr A,\mathscr O_Y),
 \qquad
 \Gamma(Y'_{\mathscr J}/Y_{\mathscr J})
 \simeq
 \operatorname{Hom}_{\mathscr O_{Y_{\mathscr J}}\text{-}\mathrm{alg}}
       (\mathscr A_{\mathscr J},\mathscr O_{Y_{\mathscr J}}).
\]

% Source page: Exp3 component-local p. 9; combined-reader p. 119.
在这些同构下, 截面 \(u=(\delta\circ g,\operatorname{id})\)
(分别地, \(u_{\mathscr J}\)) 对应于 \(\varepsilon\)
(分别地, \(\varepsilon_{\mathscr J}\)). 因此, \(P(g_{\mathscr J})\)
与模去理想后化为 \(\varepsilon_{\mathscr J}\) 的代数态射
\(\mathscr A\to\mathscr O_Y\) 一一对应; 在此双射下,
\(g\) 对应于 \(\varepsilon\).

令 \(M=g^*(\Omega^1_{X/S})\). 则
\(\operatorname{Hom}_{\mathscr O_Y\text{-}\mathrm{alg}}
 (\mathscr A,\mathscr O_Y)\) 可等同于像为平方零理想的
\(\mathscr O_Y\)-态射 \(\psi:M\to\mathscr O_Y\) 组成的集合.
我们关注其中诱导零态射
\(M\to\mathscr O_{Y_{\mathscr J}}=\mathscr O_Y/\mathscr J\)
的那些态射, 即将 \(M\) 映到 \(\mathscr J\) 内的态射.
反之, 由于 \(\mathscr J^2=0\), 每个 \(\mathscr O_Y\)-态射
\(\phi:M\to\mathscr J\) 都来自唯一一个模去理想后化为
\(\varepsilon_{\mathscr J}\) 的代数态射 \(\mathscr A\to\mathscr O_Y\).
最后,
\[
 \operatorname{Hom}_{\mathscr O_Y}
   \bigl(g^*(\Omega^1_{X/S}),\mathscr J\bigr)
 =
 \operatorname{Hom}_{\mathscr O_{Y_{\mathscr J}}}
   \bigl(g_{\mathscr J}^*(\Omega^1_{X/S}),\mathscr J\bigr)
\]
这是因为 \(\mathscr J^2=0\), 如同
\SGARefLink{sga3:iii:theorem:0-1-8:item:b}{(b)} 的证明.
从而得到一个双射
\[
 P(g_{\mathscr J})\simeq
 \operatorname{Hom}_{\mathscr O_{Y_{\mathscr J}}}
 \bigl(g_{\mathscr J}^*(\Omega^1_{X/S}),\mathscr J\bigr)
\]
在此双射下, \(g\) 对应于零态射.

对每个
\[
 m\in
 \operatorname{Hom}_{\mathscr O_{Y_{\mathscr J}}}
 \bigl(g_{\mathscr J}^*(\Omega^1_{X/S}),\mathscr J\bigr),
\]
用 \(m\cdot g\) 表示上述双射中与 \(g\) 和 \(m\) 关联的
\(P(g_{\mathscr J})\) 的元素. 我们已经看到 \(0\cdot g=g\);
还需证明
\begin{equation}
 m'\cdot(m\cdot g)=(m+m')\cdot g.
 \tag{0.1.8\,(*)}\label{sga3:iii:theorem:0-1-8:equation:0-1-8-star}
\end{equation}
这一点可局部验证.\footnote{编者注: 前面的论证几乎逐字复述
SGA~1, III~5.1; 以下作了展开.}
事实上, 上述两个态射 \(Y\to X\) 在底拓扑空间上诱导的连续映射
都与 \(g\) 相同. 因此只须验证: 对 \(X\) 的每个在 \(S\) 的仿射开集
\(\operatorname{Spec}(\Lambda)\) 上的仿射开集
\(U=\operatorname{Spec}(A)\), 以及 \(g^{-1}(U)\) 的每个仿射开集
\(V=\operatorname{Spec}(B)\), 它们诱导同一个 \(\Lambda\)-代数态射
\(A\to B\).

令 \(J=\Gamma(V,\mathscr J)\), 并令
\(\varphi,\psi,\eta:A\to B\) 分别为 \(g\), \(m\cdot g\)
和 \(m'\cdot(m\cdot g)\) 诱导的态射. 它们模去 \(J\) 后相同.
有唯一的表达式
\[
 \psi=\varphi+D,
 \qquad
 \eta=\psi+D',
\]
其中 \(D\) 属于
\[
 \operatorname{Der}_{\varphi}(A,J)
 =\{\delta\in\operatorname{Hom}_{\Lambda}(A,J)
   \mid \delta(ab)=\varphi(a)\delta(b)+\varphi(b)\delta(a)\},
\]
而 \(D'\) 属于 \(\operatorname{Der}_{\psi}(A,J)\).
但由于 \(J^2=0\), 有 \(\operatorname{Der}_{\varphi}(A,J)
 =\operatorname{Der}_{\psi}(A,J)\), 且两者都可等同于
\[
 \operatorname{Hom}_{B/J}
 \bigl(\Omega^1_{A/\Lambda}\otimes_AB/J,J\bigr).
\]
在此等同下, \(D\) 对应于 \(m\), \(D'\) 对应于 \(m'\).
因此 \(\eta=\varphi+D+D'\), 而 \(D+D'\) 对应于 \(m+m'\),
这就证明了 (\SGARefLink{sga3:iii:theorem:0-1-8}{0.1.8}\,(*)).

\medskip
\noindent\SGAInlineRefTarget{sga3:iii:corollary:0-1-9}\textbf{推论 0.1.9.}
\footnote{编者注: 此推论是在 SGA~1, III~5.1 的基础上补入的;
它证明了下面的 \SGARefLink{sga3:iii:proposition:0-2}{命题~0.2} 的 (i).}
\textit{设 \(X\) 是 \(S\)-概形, 沿用
\SGARefLink{sga3:iii:definition:0-1-5}{0.1.5} 的记号.
则 \(X\) 带有交换 \(X^+\)-群 \(L_X\) 的左作用, 使 \(X\)
在 \(X^+\) 上关于 \(L_X\) 为形式主齐性的.
换言之, 有一个 \(X^+\)-函子同构}
\[
 L_X\times_{X^+}X\xrightarrow{\sim}X\times_{X^+}X,
 \qquad
 (m,x)\longmapsto(x,m\cdot x).
\]

% Source page: Exp3 component-local p. 10; combined-reader p. 120.
\noindent\textit{证明.}
用 \(i_0:X_0\hookrightarrow X\) 表示该浸入.
由于 \(X_0=X\times_SS_0\), 首先有
\[
 i_0^*(\Omega^1_{X/S})\simeq\Omega^1_{X_0/S_0}
\]
(EGA~IV, 16.4.5).

设 \(Y\) 是一个由 \(S\)-态射 \(g_{\mathscr J}:Y_{\mathscr J}\to X\)
给定的 \(X^+\)-概形, \(g_0:Y_0\to X_0\) 由基变换得到.
根据 \SGARefLink{sga3:iii:theorem:0-1-8}{0.1.8}, 若
\(\operatorname{Hom}_{X^+}(Y,X)\) 非空, 则它在
\[
 \operatorname{Hom}_{\mathscr O_{Y_0}}
 \bigl(g_0^*i_0^*(\Omega^1_{X/S}),\mathscr J\mathscr O_Y\bigr),
\]
下为主齐性集, 后者可等同于
\[
 \operatorname{Hom}_{\mathscr O_{Y_0}}
 \bigl(g_0^*(\Omega^1_{X_0/S_0}),\mathscr J\mathscr O_Y\bigr)
 =L_X(Y).
\]
于是有一个双射
\[
 L_X(Y)\times\operatorname{Hom}_{X^+}(Y,X)
 \xrightarrow{\sim}
 \operatorname{Hom}_{X^+}(Y,X\times_{X^+}X),
 \qquad
 (m,g)\longmapsto(g,m\cdot g).
\]
我们来证明它关于 \(Y\) 具有函子性. 设 \(f:Z\to Y\) 是 \(S\)-概形态射.
只须证明下图交换.

\SGAThreeNativeDiagram[0.88\textwidth]
  {figures/exp3/diagram_011_exp3_p010_corollary019_functoriality.tex}
  {主齐性双射关于 \(Y\) 的函子性.}
  {fig:exp3-corollary-019-functoriality}
  {Polo--Gille Exp3-14oct24.pdf, 组件第 10 页; 合卷第 120 页}

若 \(\operatorname{Hom}_{X^+}(Y,X)=\emptyset\), 则无须证明.
因此只须证明: 对每个扩张 \(g_{\mathscr J}\) 的 \(S\)-态射
\(g:Y\to X\) 和每个 \(m\in L_X(Y)\), 有
\begin{equation}
 (m\cdot g)\circ f=L_X(f)(m)\cdot(g\circ f).
 \tag{0.1.9\,(*)}\label{sga3:iii:corollary:0-1-9:equation:0-1-9-star}
\end{equation}
这两个 \(S\)-态射 \(Z\to X\) 在 \(Z_{\mathscr J}\) 上都与
\(g_{\mathscr J}\circ f_{\mathscr J}\) 相同; 特别地,
它们在底拓扑空间上诱导的连续映射都与 \(g\circ f\) 相同.
因此只须在局部环上验证此等式. 设 \(z\in Z\), 令 \(y=f(z)\),
\(x=g(y)\), 并用 \(A,B,C\) 分别表示局部环
\(\mathscr O_{X,x},\mathscr O_{Y,y},\mathscr O_{Z,z}\).
设 \(s\) 是 \(x\) 在 \(S\) 中的像, 令
\(\Lambda=\mathscr O_{S,s}\), 用 \(J\) 和 \(I\) 表示
\(\Lambda\) 中分别与 \(\mathscr J\) 和 \(\mathscr I\) 对应的理想,
并令 \(\varphi,\psi:A\to B\) 及 \(\theta:B\to C\)
分别为 \(g\), \(m\cdot g\), \(f\) 诱导的 \(\Lambda\)-代数态射.
则 \(m\) 诱导下列集合的一个元素 \(D\):
\[
 \operatorname{Hom}_{B/IB}
 \bigl(\Omega^1_{A/\Lambda}\otimes_AB/IB,JB\bigr)
 =\operatorname{Der}_{\Lambda}(A,JB),
\]
并且 \(\psi=\varphi+D\). 因而 \((m\cdot g)\circ f\) 对应于
\[
 \theta\circ\psi=\theta\circ\varphi+\theta\circ D.
\]
根据 \SGARefLink{sga3:iii:definition:0-1-5}{0.1.5}, \(\theta\circ D\)
是 \(L_X(f)(m)\) 在
\[
 \operatorname{Hom}_{C/IC}
 \bigl(\Omega^1_{A/\Lambda}\otimes_AC/IC,JC\bigr)
 =\operatorname{Der}_{\Lambda}(A,JC).
\]
中的像. 因此, \(\theta\circ\varphi+\theta\circ D\) 是
\(L_X(f)(m)\cdot(g\circ f)\) 的像.
这就证明了 (\SGARefLink{sga3:iii:corollary:0-1-9}{0.1.9}\,(*)).

\medskip
\noindent\SGAInlineRefTarget{sga3:iii:corollary:0-1-10}\textbf{推论 0.1.10.}
\footnote{编者注: 这里明确表述了 \(L_X\) 关于 \(X\) 的函子性,
尤其是下面的 \SGARefLink{sga3:iii:corollary:0-1-10:item:b}{(b)}.}
\begin{enumerate}
\item[\SGAInlineRefTarget{sga3:iii:corollary:0-1-10:item:a}(a)]
  \textit{函子 \(L_X\) 函子性地依赖于 \(X\).
  对每个 \(S\)-态射 \(f:X\to W\), 有一个 \(S\)-态射
  \(L_f:L_X\to L_W\), 它是 \(f^+\) 上的交换群态射;
  也就是说, 下图交换,}

  % Source page: Exp3 component-local p. 11; combined-reader p. 121.
  \SGAThreeNativeDiagram[0.46\textwidth]
    {figures/exp3/diagram_012_exp3_p011_corollary0110_lf_square.tex}
    {\(f^+\) 上的态射 \(L_f\).}
    {fig:exp3-corollary-0110-lf-square}
    {Polo--Gille Exp3-14oct24.pdf, 组件第 11 页; 合卷第 121 页}

  \textit{并且, 对每个 \(Y\to X^+\), 映射}
  \[
   L_f(Y):\operatorname{Hom}_{X^+}(Y,L_X)
   \longrightarrow\operatorname{Hom}_{W^+}(Y,L_W),
  \]
  \textit{是交换群同态, 其中 \(Y\) 经由 \(f^+\) 在 \(W^+\) 上.}

\item[\SGAInlineRefTarget{sga3:iii:corollary:0-1-10:item:b}(b)]
  \textit{此外, 下图交换.}

  \SGAThreeNativeDiagram[0.66\textwidth]
    {figures/exp3/diagram_013_exp3_p011_corollary0110_action_square.tex}
    {\(L_f\) 与形式主齐性作用的相容性.}
    {fig:exp3-corollary-0110-action-square}
    {Polo--Gille Exp3-14oct24.pdf, 组件第 11 页; 合卷第 121 页}
\end{enumerate}

\noindent\textit{证明.}
对于 \SGARefLink{sga3:iii:corollary:0-1-10:item:a}{(a)}, \(L_f\)
由 \(\mathscr O_{X_0}\)-模态射
\[
 f_{X_0/W_0/S_0}:f_0^*(\Omega^1_{W_0/S_0})
 \longrightarrow\Omega^1_{X_0/S_0}
\]
诱导 (参见 \SGARefLink{sga3:iii:paragraph:0-1-7:item:b}{0.1.7(b)}).
对每个由 \(S\)-态射 \(g_{\mathscr J}:Y_{\mathscr J}\to X\)
给定的 \(X^+\)-概形 \(Y\), 有下列关于 \(Y\) 具有函子性的交换图.

\SGAThreeNativeDiagram[0.96\textwidth]
  {figures/exp3/diagram_014_exp3_p011_corollary0110_hom_square.tex}
  {在微分模描述下的同态 \(L_f(Y)\).}
  {fig:exp3-corollary-0110-hom-square}
  {Polo--Gille Exp3-14oct24.pdf, 组件第 11 页; 合卷第 121 页}

这里 \(L_f(Y)\) 是映射
\[
 \psi\longmapsto
 \psi\circ g_0^*(f_{X_0/W_0/S_0}),
\]
它确实是交换群同态.\footnote{编者注: 注意,
\(L_f\) 仅依赖于 \(f_0:X_0\to W_0\).}

现在证明 \SGARefLink{sga3:iii:corollary:0-1-10:item:b}{(b)}.
设 \(Y\) 是 \(X^+\)-概形. 若
\(\operatorname{Hom}_{X^+}(Y,X)=\emptyset\), 则无须证明.
因而设 \(g\in\operatorname{Hom}_{X^+}(Y,X)\). 必须证明,
对每个 \(m\in L_X(Y)\), 有
\begin{equation}
 f\circ(m\cdot g)=L_f(Y)(m)\cdot(f\circ g).
 \tag{0.1.10\,(*)}\label{sga3:iii:corollary:0-1-10:equation:0-1-10-star}
\end{equation}
固定 \(g\) 后, \(\operatorname{Hom}_X(Y,X\times_{X^+}X)\)
是 \(\operatorname{Hom}_X(Y,\Delta^{(1)}_{X/S})\) 的子集, 且
\[
\begin{aligned}
 L_X(Y)
 &=\operatorname{Hom}_{\mathscr O_{Y_0}}
   \bigl(g_0^*(\Omega^1_{X_0/S_0}),\mathscr J\mathscr O_Y\bigr)\\
 &=\operatorname{Hom}_{\mathscr O_Y}
   \bigl(g^*(\Omega^1_{X/S}),\mathscr J\mathscr O_Y\bigr),
\end{aligned}
\]
它是 \(\widetilde L_X(Y)\) 的子集; 参见
\SGARefLink{sga3:iii:paragraph:0-1-7}{0.1.7}. 此外, \(L_f(Y)\)
是 \(\widetilde L_f(Y)\) 在 \(L_X(Y)\) 上的限制.
最后, 双射
\[
 L_X(Y)\xrightarrow{\sim}
 \operatorname{Hom}_X(Y,X\times_{X^+}X),
 \qquad
 m\longmapsto(g,m\cdot g),
\]
是 \SGARefLink{sga3:iii:remark:0-1-7-1}{0.1.7.1} 中考虑的双射
\[
 \operatorname{Hom}_X(Y,\Delta^{(1)}_{X/S})
 \xrightarrow{\sim}\{g\}\times\widetilde L_X(Y)
\]
在 \(L_X(Y)\subset\widetilde L_X(Y)\) 上的限制的逆.
因此, 等式 (\SGARefLink{sga3:iii:corollary:0-1-10}{0.1.10}\,(*))
由 (\SGARefLink{sga3:iii:paragraph:0-1-7}{0.1.7}\,(*)) 得到.
% Source page: Exp3 component-local p. 12; combined-reader p. 122.
事实上, 令 \(g'\) 为由 \((g,m\cdot g)\) 定义的 \(X\)-态射
\(Y\to\Delta^{(1)}_{X/S}\). 在 \(\widetilde L_W(Y)\) 中,
与 \((f\circ g,f\circ g')\) 对应的元素是
\(\widetilde L_f(Y)(m)=L_f(Y)(m)\). 因此
\[
 L_f(m)\cdot(f\circ g)=f\circ(m\cdot g).
\]

% Continue with Lemma 0.1.11.
% Translation statistics: source_lines=591; output_lines=569; inline_math=311; display_math=38; equation_environments=3; footnotes=7.
% Review: complete assigned component; pending_TODO=0; math-token multiset, equation blocks, targets, references and environments checked.
